\begin{table}[H]
  \centering
  \caption{课程对应}
  \label{tab:concept-mapping}
  \begin{tabularx}{\linewidth}{lY}
    \toprule
    知识点 & 项目体现 \\
    \midrule
    量子态 & 线路中的门序列描述量子态演化，路由输出必须保持逻辑语义等价。\\
    张量积 & 多量子比特寄存器位于张量积空间，双比特门作用在两个子系统上。\\
    酉门 & CNOT、CZ、SWAP 都是酉门；额外 SWAP 改变承载位置，不改变理想演化。\\
    纠缠 & GHZ、QFT 等样例包含密集双比特交互，对硬件拓扑更敏感。\\
    拓扑 & 芯片耦合关系被建模为图，只有图边上的物理量子比特能直接执行双比特门。\\
    噪声 & 额外 SWAP 和更深线路会增加噪声暴露时间，减少交换门具有物理意义。\\
    优化 & 映射选择、候选 SWAP、搜索剪枝构成受约束的组合优化过程。\\
    验证 & 轨迹复放逐步检查门操作和映射变化，确认输出线路满足拓扑约束。\\
    \bottomrule
  \end{tabularx}
\end{table}
